% =====================================================================
% Patches for EvolutionaryGenerator_v2.1_Formalization_12.tex -> Revision 13
%
% Every block below is paste-ready and uses only macros and labels that already
% exist in the document. Each is preceded by its insertion point.
%
% PATCH 11 reports a campaign that is still running: 858 instances of a design
% that calls for 14400, with n=2000 set aside. It supersedes an earlier pilot that
% was confounded with the mutation regime (PATCH 4) and is reported nowhere. Its
% control experiment lives in tools/TargetMeasure.java and runs/target_measure.csv.
% Re-run
%   python3 tools/characterise.py --out runs/pisinger_main --skip-generate --lambda-m 11
% to refresh the numbers before the revision is frozen.
% =====================================================================


% ---------------------------------------------------------------------
% PATCH 1. Abstract. Append at the end of the abstract, after the Revision 12
% sentence ("...and a pilot campaign on the golden grid.").
% ---------------------------------------------------------------------

Revision~13 turns the mutation axis from a parameter into a measured result. It records that
the four configurations of the golden grid, which declare one mutation rate throughout, in
fact span a factor of $19$ in the expected number of flips per offspring
(Section~\ref{sec:mutgrid}); that holding that number fixed instead of the rate changes the
fraction of well-labelled runs from $62\%$ to $91\%$ at one configuration; and that the
eighteen mislabelled files of the reference set have two different causes, of which only one
is a property of the objective (Remark~\ref{rem:twocauses}). That second point bears on the
motivation for $\Atwo$, which is stated for all eighteen and holds for ten of them. It also
records that the golden-file kit, cited throughout as \textbf{[kit]}, is now versioned, and
that running it reproduces Remark~\ref{rem:tiesmeasured} entry by entry from an independent
implementation of the objective. Section~\ref{sec:attainability} then carries the
characterisation to the scales of Pisinger (2005), at stated rather than tuned engine settings:
the eight types split into four the generator delivers at every $n$ and four it loses as $n$
grows, the coefficient scale is almost irrelevant as Proposition~\ref{prop:scale} predicts, and
the split is \emph{not} explained by the contraction of the target sets --- two sets of equal
rarity under $\Utheta$ are attained $23/23$ and $0/24$ at the same $n$ and budget
(Remark~\ref{rem:notmeasure}). It also records that the emitted correlation $\hat\rho(w,p)$,
which the aggregate reports as zero, separates by instance type once $n$ is large enough for
the sampling noise to fall below the effect. That last result qualifies
Remark~\ref{rem:a3boundary}: the encoding makes correlation expensive rather than unreachable,
and the direction it takes is the one the correlation gene of Section~\ref{sec:a3} finds when
given a direct lever.


% ---------------------------------------------------------------------
% PATCH 2. Glossary. Insert as a new row immediately after the row for
%   "$p_c,p_m$ & crossover rate, per-bit mutation rate & Sec.~\ref{sec:engine} \\"
% ---------------------------------------------------------------------

$\lambda_m=p_mL$ & expected number of bit flips per offspring; the effective mutation axis, on which a level means the same thing at two values of $L$ & Rem.~\ref{rem:mutsemantics}, Sec.~\ref{sec:mutgrid} \\
$\alpha=c/\sum_iw_i$ & capacity ratio: the fraction of the total weight the knapsack holds; fixed by $\theta$ in $\Aone$, $\Astar$, $\Atwo$, and a gene in $A_3^{c}$ & Rem.~\ref{rem:theta3alpha}, Def.~\ref{def:a3genotype} \\


% ---------------------------------------------------------------------
% PATCH 3. Section 7.4. Insert immediately after the end of
%   \begin{remark}[Measured on the reference set]\label{rem:tiesmeasured} ... \end{remark}
% and before \begin{remark}[Why this cannot be a change inside $\Astar$].
% ---------------------------------------------------------------------

\begin{remark}[The eighteen have two causes, not one]\label{rem:twocauses}
The failures above split by cause, and only one of the two is a property of the objective. At
$\theta_2$ the shipped $p_m=0.1$ is $\lambda_m=55$ flips per offspring
(Section~\ref{sec:mutgrid}), inside the regime in which the engine stops searching. Re-running
$\theta_2$ on the same ten seeds of the reference set, the same budget $\Emax$ and the same
everything else, at $p_m=0.02$, that is $\lambda_m=11$ \evid{V3}:

\begin{center}\small
\begin{tabular}{@{}lccccc@{}}
\toprule
type at $\theta_2$ & $p_m$ & $\lambda_m$ & well labelled & median $f_\tau$ \\
\midrule
\code{DEFAULT\_EASY}             & $0.1$  & $55$ & $2/10$           & $-10$ \\
\code{DEFAULT\_EASY}             & $0.02$ & $11$ & $\mathbf{10/10}$ & $+29$ \\
\code{MAX\_PROFIT\_WEIGHT\_HARD} & $0.1$  & $55$ & $0/10$           & $-42$ \\
\code{MAX\_PROFIT\_WEIGHT\_HARD} & $0.02$ & $11$ & $0/10$           & $-11$ \\
\bottomrule
\end{tabular}
\end{center}

The $2/10$ of the first row is the $q_{\theta,\tau}=0.2$ of
Proposition~\ref{prop:conditional}, recomputed here as a check that the two runs are the same
experiment. The conclusion is that the two failing types are not the same phenomenon:

\begin{itemize}[leftmargin=1.6em, itemsep=2pt]
\item The eight \code{DEFAULT\_EASY} failures are an artefact of the mutation regime. They
disappear at the standard budget of $10^{5}$ evaluations once $\lambda_m$ is brought back to
the value the shipped defaults give at $\theta_0$. Remark~\ref{rem:a2split} reports $\Atwo$
reaching $8/8$ on the same cell by spending up to $661\,220$ evaluations; lowering the
mutation rate reaches $10/10$ at $\Emax$.
\item The ten \code{MAX\_PROFIT\_WEIGHT\_HARD} failures are not. The deficit improves from
$-42$ to $-11$ and never turns positive, consistent with $\Utheta(d_{\MAXPW}>0)<10^{-5}$ at
$\theta_2$ (Section~\ref{sec:why}) and with $\Atwo$ also failing $0/10$ at ten times the
budget.
\end{itemize}

The attainability map of Section~\ref{sec:why} therefore has to distinguish a cell the search
cannot reach from a cell the shipped engine settings cannot reach. Only the second kind is a
statement about the geometry of the portfolio.
\end{remark}


% ---------------------------------------------------------------------
% PATCH 4. Section 9. Insert as a new subsection at the very end of the section,
% after \begin{kept}{the published tuning} ... \end{kept} and before
%   % =====================================================================
%   \section{Randomness architecture}
% ---------------------------------------------------------------------

\subsection{The mutation axis on the grid itself}\label{sec:mutgrid}

\begin{whatis}{why one number replaces another}
Remark~\ref{rem:mutsemantics} observes that \code{mutate} flips every bit independently, so
what an offspring actually receives is $\lambda_m=p_mL$ flips, not $p_m$. Section~\ref{sec:lhsverdict}
then recommends parametrising the mutation axis by $\lambda_m$. This subsection applies that
recommendation to the document's own grid, where it has a consequence that was not drawn: the
four configurations are not comparable to each other.
\end{whatis}

The golden grid declares a single mutation rate, $p_m=0.1$, at all four configurations. In
flips per offspring they are four different regimes \evid{V3}:

\begin{center}\small
\begin{tabular}{@{}lrrrc@{}}
\toprule
 & $n$ & $L$ & $\lambda_m$ at $p_m=0.1$ & well labelled \\
\midrule
$\theta_0$ & $10$  & $110$  & $11$  & $80/80$ \\
$\theta_1$ & $12$  & $120$  & $12$  & $80/80$ \\
$\theta_2$ & $50$  & $550$  & $55$  & $62/80$ \\
$\theta_3$ & $100$ & $2100$ & $210$ & $32/32$ \\
\bottomrule
\end{tabular}
\end{center}

A factor of $19$ separates the extremes. Whatever is observed to change across the grid is
therefore confounded with a change of mutation regime, and a level of the nominal axis does
not name a level of the effective one --- the same collapse Proposition~\ref{prop:lhspush}
describes for the coefficient axis, here on the axis the 2019 study did tune.

$\lambda_m$ does not by itself order the outcome. $\theta_3$ carries the largest value and is
labelled perfectly; its capacity ratio $\alpha=0.007$ leaves a nearly empty knapsack on which
the four heuristics separate easily (Remark~\ref{rem:theta3alpha}), so the task is easy there
whatever the mutation. Isolating the axis needs the configuration held fixed. At $n=100$,
$\Wmax=\Pmax=1024$ (so $\bw=\bp=10$, $L=2000$), $c=25\,625$, the eight labelled types, four
instances each, $p_c=0.8$, $k=3$, $\mu=10$, budget $\Emax$ \evid{V3}:

\begin{center}\small
\begin{tabular}{@{}rrcc@{}}
\toprule
$\lambda_m$ & $p_m$ & well labelled & median $f_\tau$ \\
\midrule
$11$  & $5.5\times10^{-3}$ & $\mathbf{29/32}$ & $\mathbf{6087}$ \\
$50$  & $2.5\times10^{-2}$ & $20/32$ & $2342$ \\
$200$ & $1.0\times10^{-1}$ & $20/32$ & $218$ \\
\bottomrule
\end{tabular}
\end{center}

The three values are not arbitrary: $11$ is what the shipped defaults give at $\theta_0$,
$200$ is what keeping $p_m=0.1$ gives at this configuration, and $50$ is between them and
close to $\theta_2$'s $55$. \code{DEFAULT\_EASY} and \code{MIN\_WEIGHT\_EASY} go from $0/4$ to
$4/4$; \code{MAX\_PROFIT\_WEIGHT\_HARD} from $0/4$ to $1/4$; the other five types are $4/4$
throughout. The median fitness moves by a factor of $28$. This is the controlled form of the
observation in Section~\ref{sec:why} that the engine loses to best-of-sample in the thin
cells: it does so at $\lambda_m=55$, not at $\lambda_m=11$.

\begin{kept}{the shipped mutation rate}
$p_m$ remains the parameter the code exposes, and $0.1$ remains what the grid and the
reference set were recorded with. Changing it changes every emitted file, so it cannot be
changed inside $\Astar$. What the thesis adopts is the reporting convention: every campaign
states $\lambda_m$ alongside $p_m$, and a design that varies $n$ or the coefficient widths
holds $\lambda_m$ fixed rather than $p_m$.
\end{kept}

\begin{remark}[What is not settled]\label{rem:lambdaopen}
Holding $p_m$ fixed and holding $\lambda_m$ fixed preserve different things: the first
preserves the fraction of flipped bits, and with it the expected flips per item block
$p_m\ell$; the second preserves the absolute number, and with it the expected number of items
a mutation touches. They coincide only at fixed $\ell$ and $n$. The measurement above shows
that fewer flips beat more at one configuration; it does not show that $\lambda_m$ is the
invariant that makes two configurations comparable, and no argument here decides it. Settling
it needs $\lambda_m$ swept as an axis of the design, which the campaign tool supports and
which costs a fraction of the characterisation runs.
\end{remark}


% ---------------------------------------------------------------------
% PATCH 5. Section 6, after Proposition~\ref{prop:dyadic} and its Code note,
% before the nominal/effective figure.
% ---------------------------------------------------------------------

\begin{remark}[The dyadic collapse at Pisinger's scales]\label{rem:pisingerR}
The collapse is not only a statement about small grids. Pisinger (2005) indexes his instance
classes by a coefficient range $R\in\{10^{3},10^{4},10^{5}\}$; the generator can realise only
powers of two, so those become \evid{V2}
\begin{center}\small
\begin{tabular}{@{}lccc@{}}
\toprule
$R$ & $\bw=\lceil\log_2R\rceil$ & coefficients reachable & departure from $R$ \\
\midrule
$10^{3}$ & $10$ & $1..1024$   & $+2\%$ \\
$10^{4}$ & $14$ & $1..16384$  & $+64\%$ \\
$10^{5}$ & $17$ & $1..131072$ & $+31\%$ \\
\bottomrule
\end{tabular}
\end{center}
Only $R=10^{3}$ has a faithful image; $10^{4}$ has none, since the neighbouring level
$2^{13}=8192$ is $18\%$ below. Within a level nothing moves at all: every $\Wmax\in[513,1024]$
is the same generator, and only the file name records which was asked for. A campaign at
Pisinger's scales should therefore state the realised $2^{\bw}$, not the requested $R$.
\end{remark}


% ---------------------------------------------------------------------
% PATCH 6. Section 14, immediately after the \begin{whatis}{the golden file}
% box and before "The grid, the seeds and the engine parameters are the literals...".
% ---------------------------------------------------------------------

\begin{remark}[$\theta_3$ has Pisinger's scale but not its regime]\label{rem:theta3alpha}
$\theta_3$ is named for the coefficient magnitude of Pisinger (2005) and matches it
(Remark~\ref{rem:pisingerR}). Its capacity does not. With $c=400$ and
$\sum_iw_i\approx5.4\times10^{4}$ the ratio is $\alpha=0.007$, and on a reference file the
four heuristics pack between $4$ and $7$ of the $100$ items \evid{V3}, where Pisinger draws
the capacity as a fraction of the total weight, commonly one half. Results at $\theta_3$
therefore describe a nearly empty knapsack. That is also why $\theta_3$ is labelled $32/32$
despite carrying the largest $\lambda_m$ of the grid (Section~\ref{sec:mutgrid}): with so few
items packed the four heuristics separate easily, and a strict winner is cheap. The
configuration is kept as recorded --- it is part of the golden file --- but no claim about
the Pisinger regime rests on it.
\end{remark}


% ---------------------------------------------------------------------
% PATCH 7. Section 14, replacing the paragraph that begins
%   "Negative control, recorded: correcting \code{Statistical.max} to $-\infty$ ..."
% (keep that paragraph, and add the two below after it).
% ---------------------------------------------------------------------

\paragraph{The kit, as versioned.} The kit cited as \textbf{[kit]} throughout is
\code{tests/golden/} of the repository: \code{GoldenGrid.java}, \code{golden.py} with the
three subcommands \code{record}, \code{verify} and \code{label}, and \code{reference/} holding
the $544$ files, their \code{MANIFEST.sha256} and \code{NOT\_WELL\_LABELLED.txt}. The manifest
was recorded from $\Astar$ rather than from $\Aone$, whose tree no longer exists as a working
copy; the two are output-identical on the grid, which is what
Section~\ref{sec:obligations} establishes, so the manifest is the oracle of $\Aone$ as stated.
\code{golden.py verify} against $\Astar$ reports $544/544$ \evid{V3}, which extends to the
whole grid an identity previously checked on a subset. Against $\Azero$ it fails on most
files, by Proposition~\ref{prop:crossover}; that is the intended behaviour of the instrument,
not a regression.

\paragraph{Remark~\ref{rem:tiesmeasured} as a reproducible result.} \code{golden.py label}
recomputes the objective from the emitted files with an independent implementation of the four
priority keys, the single first-fit pass and the guarded extrema, and never calls the
generator. It reproduces the table of Remark~\ref{rem:tiesmeasured} entry by entry \evid{V3}:
$18$ of $272$ labelled files not well labelled ($6.6\%$); eight \code{DEFAULT\_EASY} and ten
\code{MAX\_PROFIT\_WEIGHT\_HARD}, all at $\theta_2$; minima and medians $-20$ / $24$ and
$-54$ / $13$; and $\ge13$ and $\ge52$ for the two groups of three that never fail. That the
numbers come out of a second implementation is the point: it checks the recorded grid against
the document rather than against itself.


% ---------------------------------------------------------------------
% PATCH 8. Section 15, immediately after the \begin{whatis}{the correction} box.
% ---------------------------------------------------------------------

One qualification belongs here rather than later. Of the eighteen files,
Remark~\ref{rem:twocauses} attributes eight to the mutation regime and ten to the objective:
the eight \code{DEFAULT\_EASY} failures at $\theta_2$ are repaired at the standard budget by
returning $\lambda_m$ to the value the shipped defaults give at $\theta_0$, whereas the ten
\code{MAX\_PROFIT\_WEIGHT\_HARD} failures survive both that change and, by
Remark~\ref{rem:a2split}, ten times the budget. $\Atwo$ is therefore motivated on the ten, not
on the eighteen. It remains the right construction for them --- the extended search is what
distinguishes a cell the engine has not yet reached from one it cannot reach --- but the
figure to quote for what it repairs is not the whole of the mislabelled tail, and a reader
comparing $\Atwo$ against a cheaper remedy should be told that a cheaper remedy exists for
about half of it.


% ---------------------------------------------------------------------
% PATCH 9. Table 2 (tab:changes), row 5. Optional. Replace the "Licensed by"
% cell of change 5 so the parallelism claim is scoped to what was verified.
% ---------------------------------------------------------------------

by construction plus golden file; removes the shared static encoding state that prevents two configurations from coexisting in one JVM (thread safety of the rest is not claimed)


% ---------------------------------------------------------------------
% PATCH 10. Section 17, item 3 of the summary. Append one sentence.
% ---------------------------------------------------------------------

The mutation rate enters as $\lambda_m=p_mL$ flips per offspring, which is the level that must
be held fixed across configurations; at the shipped $p_m=0.1$ the golden grid itself spans
$\lambda_m$ from $11$ to $210$ (Section~\ref{sec:mutgrid}).


% ---------------------------------------------------------------------
% PATCH 11. Section 7. Insert as a new subsection at the end of the section,
% after \subsection{Selection hardness and solver hardness}\label{sec:hardness}
% and before the  % ===== \section{The evolutionary engine}  line.
% ---------------------------------------------------------------------

\subsection{Attainability away from the golden grid}\label{sec:attainability}

\begin{whatis}{what this campaign asks and at which settings}
Remark~\ref{rem:tiesmeasured} measures which labels the generator delivers on the four
configurations of the golden file, where $n\le100$. This section asks the same question at the
scales Pisinger (2005) uses. It is a characterisation of $\Gtt$ \emph{at one engine setting},
not of the generator at its best: following the red box of Section~\ref{sec:lhsverdict} the
engine is not re-tuned, and the settings are stated rather than optimised. They are
$\mu=10$ and $k=3$, the shipped defaults; $p_c=0.8$, the value the golden grid uses so that the
crossover guard is exercised; the hard-coded budget $\Emax=10^{5}$; and $\lambda_m=11$ flips
per offspring, which is what the shipped $p_m=0.1$ gives at $\theta_0$ and which
Section~\ref{sec:mutgrid} shows must be held rather than $p_m$ when $L$ varies. Only the last
departs from what the code exposes, and it departs in the direction the document itself
recommends.
\end{whatis}

The design is a full factorial on the effective axes: $n\in\{50,100,200,500,1000\}$;
$\Wmax=\Pmax\in\{2^{10},2^{14}\}$, the realisable images of Pisinger's $R=10^{3}$ and
$R=10^{4}$ (Remark~\ref{rem:pisingerR}); and a capacity ratio $\alpha\in\{0.05,0.25,0.50\}$,
with $c$ derived per cell from the expected total weight, $c=\lceil\alpha\,n(2^{\bw}+1)/2\rfloor$,
rather than carried over from another configuration --- the realised $c/\sum_iw_i$ averages
$0.050$, $0.261$ and $0.554$ against the three nominal values \evid{V3}. Eight labelled types.
The seed ignores $\alpha$, so its levels are compared on shared random numbers. The figures
below are $858$ emitted instances, the design not yet complete.

\paragraph{The portfolio splits in two, and the split sharpens with $n$.}

\begin{center}\small
\begin{tabular}{@{}lccccc@{}}
\toprule
type & $n{=}50$ & $n{=}100$ & $n{=}200$ & $n{=}500$ & $n{=}1000$ \\
\midrule
\code{DEFAULT\_HARD}             & $24/24$ & $24/24$ & $24/24$ & $24/24$ & $12/12$ \\
\code{MIN\_WEIGHT\_HARD}         & $24/24$ & $24/24$ & $24/24$ & $23/23$ & $12/12$ \\
\code{MAX\_PROFIT\_HARD}         & $24/24$ & $24/24$ & $24/24$ & $23/23$ & $12/12$ \\
\code{MAX\_PROFIT\_WEIGHT\_EASY} & $24/24$ & $24/24$ & $24/24$ & $23/23$ & $12/12$ \\
\addlinespace[2pt]
\code{MIN\_WEIGHT\_EASY}         & $24/24$ & $24/24$ & $22/24$ & $8/23$  & $2/12$ \\
\code{MAX\_PROFIT\_EASY}         & $24/24$ & $24/24$ & $24/24$ & $0/23$  & $0/12$ \\
\code{DEFAULT\_EASY}             & $24/24$ & $23/24$ & $9/24$  & $0/24$  & $0/12$ \\
\code{MAX\_PROFIT\_WEIGHT\_HARD} & $13/24$ & $2/24$  & $1/24$  & $0/23$  & $0/12$ \\
\bottomrule
\end{tabular}
\end{center}

Four types are delivered at every $n$ without a single failure. The other four are lost, in
the order \MAXPW-\textsc{hard} (already at $n=50$), \DEF-\textsc{easy} ($n=200$), then
\MAXP-\textsc{easy} and \MINW-\textsc{easy} ($n=500$). Aggregated, the rate falls from
$181/192$ at $n=50$ to $50/96$ at $n=1000$. This is the attainability map of
Section~\ref{sec:why} measured on an axis the golden grid does not vary, and it is not the
mutation artefact of Remark~\ref{rem:twocauses}: $\lambda_m$ is $11$ in every cell.

\paragraph{The coefficient scale is almost irrelevant, and so is the capacity ratio.} Between
$R=10^{3}$ and $R=10^{4}$ the counts are $326/432$ and $327/426$. That is what
Proposition~\ref{prop:scale} predicts: profits enter the packings only through comparisons and
the totals linearly, so widening the range scales every margin without changing a winner. The
capacity ratio moves it little as well ($232/288$, $207/288$ and $214/282$ at $\alpha=0.05$,
$0.25$, $0.50$). What moves attainability is $n$.

\paragraph{The measure of the target set does not explain the split.} The obvious reading of
the table is that the lost types are the ones whose target set $W_\tau$ has contracted. That
reading is wrong, and the check is cheap: draw from $\Utheta$ directly and count strict winners
and losers, with no engine, no elite and no budget, as Section~\ref{sec:why} does for the grid.
Averaged over $R$ and $\alpha$, with $10^{5}$ draws at $n=50$ down to $5\times10^{3}$ at
$n=1000$ \evid{V3}:

\begin{center}\small
\begin{tabular}{@{}lccccc@{}}
\toprule
$\Utheta(W_\tau)$ & $n{=}50$ & $n{=}100$ & $n{=}200$ & $n{=}500$ & $n{=}1000$ \\
\midrule
\code{DEFAULT\_HARD}             & $0.955$ & $0.982$ & $0.996$ & $1.000$ & $1.000$ \\
\code{MAX\_PROFIT\_WEIGHT\_EASY} & $0.975$ & $1.000$ & $1.000$ & $1.000$ & $1.000$ \\
\code{MAX\_PROFIT\_HARD}         & $0.040$ & $0.018$ & $0.004$ & $<10^{-4}$ & $<2\times10^{-4}$ \\
\code{MIN\_WEIGHT\_HARD}         & $0.005$ & $<10^{-4}$ & $<4\times10^{-5}$ & $<10^{-4}$ & $<2\times10^{-4}$ \\
\addlinespace[2pt]
\code{MAX\_PROFIT\_EASY}         & $0.008$ & $<10^{-4}$ & $<4\times10^{-5}$ & $<10^{-4}$ & $<2\times10^{-4}$ \\
\code{MIN\_WEIGHT\_EASY}         & $0.001$ & $<10^{-4}$ & $<4\times10^{-5}$ & $<10^{-4}$ & $<2\times10^{-4}$ \\
\code{DEFAULT\_EASY}             & $10^{-5}$ & $<2\times10^{-5}$ & $<4\times10^{-5}$ & $<10^{-4}$ & $<2\times10^{-4}$ \\
\code{MAX\_PROFIT\_WEIGHT\_HARD} & $<10^{-5}$ & $<2\times10^{-5}$ & $<4\times10^{-5}$ & $<10^{-4}$ & $<2\times10^{-4}$ \\
\bottomrule
\end{tabular}
\end{center}

Six of the eight target sets do contract to below the resolution of the sample: under the
uniform law at $n\ge200$, \MAXPW\ is the strict winner and \DEF\ the strict loser on virtually
every instance, and the portfolio has stopped disagreeing in any other way. But the contraction
does not sort the two halves of the previous table. At $n=500$, \code{MIN\_WEIGHT\_HARD} and
\code{DEFAULT\_EASY} both have $\Utheta(W_\tau)$ below $10^{-4}$, and the engine delivers the
first $23/23$ and the second $0/24$. \code{MAX\_PROFIT\_HARD}, equally rare, is $23/23$.

\begin{remark}[What the split is not]\label{rem:notmeasure}
A vanishing $\Utheta(W_\tau)$ is compatible with perfect attainment and with none, at the same
$n$, the same budget and the same settings. So the loss of four types with $n$ is not explained
by the measure of their target sets, and the reading that ``the generator cannot reach them
because they became rare'' is not supported: it reaches two sets of the same rarity every time.
Nor is it explained by budget exhaustion alone, for the same reason. What distinguishes the two
halves is a property of the landscape the objective induces --- whether the high-margin
instances form a basin the search can climb --- and nothing measured here decides it. The
observation stands as a fact about the shipped configuration; its mechanism is open. This is
the same distinction Section~\ref{sec:why} draws when it notes that the search reaches sets
sampling would need more than $10^{5}$ draws to hit once: rarity under $\Utheta$ says nothing
on its own about reachability under the engine.
\end{remark}

\paragraph{The emitted correlation is not centred at zero; it is centred by type.} Over the
whole campaign $\hat\rho(w,p)$ averages $-0.03$, which reproduces the familiar statement that
the generator emits uncorrelated instances. The aggregate hides the effect. At $n\ge500$, where
the sampling noise of a correlation under $\Utheta$ is below $0.045$, the type-conditional
means are \evid{V3}:

\begin{center}\small
\begin{tabular}{@{}lrr@{}}
\toprule
type & mean $\hat\rho(w,p)$ & sd \\
\midrule
\code{DEFAULT\_HARD}             & $-0.454$ & $0.115$ \\
\code{MIN\_WEIGHT\_EASY}         & $-0.171$ & $0.060$ \\
\code{MAX\_PROFIT\_EASY}         & $-0.109$ & $0.129$ \\
\code{MAX\_PROFIT\_WEIGHT\_HARD} & $+0.092$ & $0.119$ \\
\code{MAX\_PROFIT\_HARD}         & $+0.130$ & $0.146$ \\
\code{MAX\_PROFIT\_WEIGHT\_EASY} & $+0.138$ & $0.047$ \\
\code{DEFAULT\_EASY}             & $+0.189$ & $0.091$ \\
\code{MIN\_WEIGHT\_HARD}         & $+0.250$ & $0.058$ \\
\bottomrule
\end{tabular}
\end{center}

Every one of these is several standard errors from zero, and \code{DEFAULT\_HARD} is about
ten. The signs are the ones Section~\ref{sec:a3} finds when the correlation gene gives the
search a direct lever: \DEF-\textsc{hard} to the bottom, \MINW-\textsc{easy} and
\MAXP-\textsc{easy} below zero, \MAXPW-\textsc{hard} above it. The gene reaches $\rho=-1$ for
\DEF-\textsc{hard}; the shipped encoding reaches $-0.45$ for the same type by paying for it
through the item bits.

\begin{remark}[What this changes in Section~\ref{sec:a3}]\label{rem:rhoqualify}
Remark~\ref{rem:a3boundary} states that the shipped generator ``has no lever on
$\hat\rho(w,p)$'' and that its corpora ``still sit at $\hat\rho\approx0$''. The first half
stands: Proposition~\ref{prop:factor} is a statement about the fitness and nothing measured
here touches it. The second half does not. It holds on the golden grid because $n\le100$ there
and the sampling noise of a correlation, $1/\sqrt{n-1}$, is $0.33$ at $\theta_0$ and $0.10$ at
$\theta_3$ --- wide enough to hide a type-conditional mean of $0.2$. At $n\ge500$ the noise
falls below $0.05$ and the separation is plain. The encoding therefore makes correlation
\emph{expensive}, not unreachable: the objective has no term in $\hat\rho$, yet the law $\Gtt$
it induces is displaced in that coordinate, and the displacement grows readable with $n$. A
corpus of these types is not a sample of uncorrelated instances, and the support-boundary row
that says results do not transfer to correlated classes should be read as a statement about
magnitude, not about direction.
\end{remark}

\begin{kept}{the engine settings, and what this campaign does not show}
The engine is not re-tuned, so every number above is conditional on $\mu=10$, $k=3$,
$p_c=0.8$, $\Emax=10^{5}$ and $\lambda_m=11$. A different budget or population would give a
different table, and the table is not evidence about the generator at its best. Two axes of
the design of Section~\ref{sec:lhsverdict} are missing: $\lambda_m$ itself, which is held
rather than swept, and the population size. The campaign is also incomplete --- $858$
instances of a design that calls for $14\,400$, and $n=2000$ set aside --- so the counts are
quoted as they stand and no distributional claim is made. None of this is a golden-file result:
the configurations are outside the grid, so the instruments of Section~\ref{sec:obligations}
do not apply, and reproducibility here means only that the seeds and settings are recorded.
\end{kept}
